\chapter{SYSTEM ARCHITECTURE AND DESIGN}

The methodology employed in this project revolves around constructing a highly modular, step-by-step data pipeline. The core engineering philosophy is to ingest unstructured, highly noisy document data, normalize it into structured JSON objects, pass those objects through a rigorous mathematical scoring engine, and render the auditable results via an interactive dashboard interface.

\section{Overall Architecture}
The system architecture operates synchronously via a lightweight Flask backend. When a recruiter uploads a batch of resumes alongside a target Job Description (JD), the system initializes an asynchronous multi-threaded processing loop to prevent interface blocking. The pipeline consists of five primary sequential modules: the Parsing Engine, the NLP Extraction Engine, the Mathematical Matching Engine, the XAI Explainer, and the Visualization Dashboard.

The architecture guarantees that every candidate document is subjected to the exact same sequence of normalization, extraction, and scoring, ensuring a standardized baseline for comparison.

\begin{figure}[H]
    \centering
    \includegraphics[width=0.92\textwidth]{project_flowchart.png}
    \caption{AI-Powered Resume Screening and Matching Architecture}
    \label{fig:project-flowchart}
\end{figure}

\section{File Parsing and OCR Fallback Pipeline}
The parsing module is explicitly designed to eliminate format bias. It receives binary file streams and attempts to extract text using the most computationally efficient method first, falling back to heavier computer-vision methods only if the initial extraction fails validation.

\begin{enumerate}
    \item \textbf{Digital Extraction:} For .docx files, the python-docx library reads the underlying XML trees. For digital .pdf files, PyPDF2 traverses the binary object streams.
    \item \textbf{OCR Fallback Validation:} If the digital extraction yields fewer than 50 characters, the system assumes the PDF is an unsearchable scanned image or a heavily protected format.
    \item \textbf{Pillow and Tesseract OCR:} Image formats (.png, .jpg) and failed PDFs are routed directly to the Tesseract OCR engine. To maximize optical recognition accuracy, the Pillow library applies a series of localized preprocessing filters:
\end{enumerate}

\begin{itemize}
    \item \textbf{Grayscale Conversion:} Reduces RGB channel noise.
    \item \textbf{Binarization:} Thresholds pixels to pure black or white, exponentially heightening text contrast against the background canvas.
    \item \textbf{Sharpening:} Applies an edge-enhancement convolution filter to clarify small or blurred fonts before character mapping occurs.
\end{itemize}

\textbf{Algorithm 1: Resume Parsing and OCR Fallback Validation}

\begin{verbatim}
FUNCTION ParseDocument(file, extension):
    IF extension == ".txt":
        return read_text(file)
    ELSE IF extension == ".pdf":
        text = PyPDF2_extract(file)
        IF length(text) < 50:
            image = convert_pdf_to_image(file)
            return OCR_Pipeline(image)
        return text
    ELSE IF extension == ".png" OR ".jpg":
        return OCR_Pipeline(file)

FUNCTION OCR_Pipeline(image):
    gray_img = apply_grayscale(image)
    binary_img = apply_binarization(gray_img)
    sharp_img = apply_sharpening(binary_img)
    return Tesseract_Extract_Text(sharp_img)
\end{verbatim}

\section{Natural Language Processing (NLP) Module}
Once a uniform, lowercase text string is generated by the parsing pipeline, it is passed to the SpaCy NLP module, utilizing the en\_core\_web\_sm model.

\begin{itemize}
    \item \textbf{Tokenization and Stop-word Removal:} Unnecessary grammatical syntax, such as ``and'', ``the'', and ``is'', is stripped to isolate high-value nouns and verbs.
    \item \textbf{Entity Matching:} The system references a predefined ontology of hard technical skills, such as ``Python'', ``AWS'', and ``Docker'', and soft skills, such as ``Leadership'' and ``Communication''.
    \item \textbf{Years of Experience (YoE) Extraction:} Utilizing complex Regular Expressions (Regex), the module scans for chronological date patterns, such as ``Aug 2018 - Present'', or explicit numeric claims, such as ``5 years of experience'', to mathematically estimate the total career tenure of the applicant.
\end{itemize}

\section{Vectorization and Similarity Computation}
While explicit entity matching handles the strict technical boundaries, the prototype employs statistical machine learning to capture broader contextual relevance. The system uses the scikit-learn TfidfVectorizer to convert both the parsed resume and the job description into dense numerical arrays.

The Term Frequency-Inverse Document Frequency (TF-IDF) mechanism ensures that rare, highly specific words, like ``Kubernetes'', carry significantly more mathematical weight than common industry jargon, like ``developer''. This prevents applicants from scoring highly simply by repeating generic terms.

The contextual similarity is then computed using the Cosine Similarity formula, mapping the angle between the two multidimensional vectors:

\begin{equation}
\text{Cosine Similarity}(A,B)=\frac{A \cdot B}{\|A\|\|B\|}
\end{equation}

where \(A\) represents the Resume TF-IDF vector, and \(B\) represents the Job Description TF-IDF vector.

\section{The Weighted Scoring Engine}
To prevent keyword stuffing from corrupting the final result, the system utilizes a capped, weighted scoring formula rather than a flat, naive average. The recruiter can adjust the scalar weights dynamically via the frontend. By default, the system assumes \(W_{tech}=0.60\), \(W_{soft}=0.10\), \(W_{yoe}=0.15\), and \(W_{context}=0.15\). Furthermore, weight normalization logic ensures that:

\begin{equation}
W_{tech}+W_{soft}+W_{yoe}+W_{context}=1.0
\end{equation}

\begin{equation}
\begin{aligned}
\text{Final Score} ={}&
(\text{Tech Score} \times W_{tech}) +
(\text{Soft Score} \times W_{soft})\\
&+(\text{YoE Score} \times W_{yoe}) +
(\text{Context Score} \times W_{context})\\
&+\text{Extra Skills Bonus}
\end{aligned}
\end{equation}

\textbf{Algorithm 2: Extra Skills Bonus Math}

To reward candidates possessing adjacent technical skills not explicitly requested in the job description, an additive bonus is applied. Crucially, this bonus is strictly capped to prevent runaway score inflation and boundary breaking.

\begin{footnotesize}
\begin{verbatim}
FUNCTION CalculateBonus(resume_skills, jd_skills, max_bonus_weight):
    extra_skills = set(resume_skills) - set(jd_skills)
    raw_bonus = length(extra_skills) * (max_bonus_weight / 10)
    final_bonus = MIN(raw_bonus, max_bonus_weight)
    return final_bonus
\end{verbatim}
\end{footnotesize}

\section{Explainable AI (XAI) Rule Engine}
The Explainable AI module maps the raw computational vectors back into a deterministic, human-readable format. It inherently classifies candidates into four confidence buckets: Strong, Moderate, Partial, and Weak.

The \texttt{Explainer.explain()} method intercepts the mathematical dictionary generated by the Matching Engine and uses a templating engine to construct a plain-English justification. This mapping specifies exact fractions, such as ``Matched 2/4 required skills'', and lists specific missing competencies, ensuring absolute HR auditability and mitigating the black-box phenomenon entirely.
